\documentclass[a4paper,11pt]{article}

\usepackage[T1]{fontenc}
\usepackage[utf8]{inputenc}
\usepackage{lmodern}
\usepackage[french]{babel}
\usepackage[margin=2.3cm]{geometry}
\usepackage{microtype}
\usepackage{booktabs}
\usepackage{tabularx}
\usepackage{xcolor}
\usepackage{listings}
\usepackage{tikz}
\usetikzlibrary{positioning,arrows.meta,fit,calc}
\usepackage[hidelinks]{hyperref}

\definecolor{codebg}{gray}{0.96}
\definecolor{accent}{RGB}{31,78,121}
\definecolor{couche}{RGB}{225,236,246}
\definecolor{entite}{RGB}{240,240,240}

\lstset{
  basicstyle=\ttfamily\small,
  backgroundcolor=\color{codebg},
  frame=single, rulecolor=\color{gray!40},
  breaklines=true, columns=fixed, basewidth=0.5em,
  keywordstyle=\color{accent}\bfseries,
  commentstyle=\color{gray}\itshape,
  literate={é}{{\'e}}1 {è}{{\`e}}1 {à}{{\`a}}1 {ç}{{\c c}}1 {ê}{{\^e}}1
           {ô}{{\^o}}1 {î}{{\^i}}1 {â}{{\^a}}1 {ù}{{\`u}}1 {É}{{\'E}}1,
}

\newcommand{\code}[1]{\texttt{#1}}

\title{\textbf{Gestion RH du laboratoire}\\[0.3em]\large Description de l'architecture}
\author{}
\date{Octobre 2026}

\begin{document}
\maketitle

\begin{abstract}
Ce document décrit l'architecture de l'application de gestion des ressources
humaines du laboratoire : pile technique, découpage en couches, modèle de
données, conventions des contrôleurs, recherche générique, et invariants à
respecter lors des évolutions. Il complète le \code{README.md} et le
\code{CLAUDE.md} du dépôt.
\end{abstract}

\tableofcontents

%======================================================================
\section{Objectif et périmètre}

L'application gère les personnes d'un laboratoire de recherche et tout ce qui
s'y rattache : structures (équipes, services), contrats, employeurs et grades,
postes occupés, locaux (campus, bâtiments, bureaux), adresses et pays,
publications.

Les choix de la version~1 sont volontairement minimalistes :
\begin{itemize}
  \item une application web monolithique rendue côté serveur, sans frontend
        séparé ni API REST ;
  \item pas d'authentification ;
  \item toute relation du modèle est navigable dans les deux sens depuis
        l'interface, comme l'exige le cahier des charges ;
  \item tout le vocabulaire métier (classes, champs, routes, templates) est en
        français, pour coller au cahier des charges.
\end{itemize}

%======================================================================
\section{Pile technique}

\begin{center}
\begin{tabularx}{\linewidth}{@{}lX@{}}
\toprule
\textbf{Composant} & \textbf{Rôle} \\
\midrule
Java 21 & Langage et environnement d'exécution \\
Spring Boot 4.1 & Socle applicatif (starters \code{webmvc}, \code{data-jpa},
  \code{thymeleaf}, \code{validation}, \code{devtools}) \\
Spring Data JPA / Hibernate & Mapping objet-relationnel, repositories,
  \code{Specification} pour la recherche \\
Thymeleaf & Rendu HTML côté serveur \\
Jakarta Bean Validation & Contraintes déclaratives (\code{@NotBlank}) \\
MariaDB 11 & Base de données, installée sur le poste (Homebrew) \\
Maven & Build (wrapper \code{mvnw} fourni) \\
\bottomrule
\end{tabularx}
\end{center}

%======================================================================
\section{Vue d'ensemble en couches}

L'application suit une architecture MVC classique en trois couches, sans
couche « service » intermédiaire : la logique métier est réduite et les
contrôleurs parlent directement aux repositories.

\begin{center}
\begin{tikzpicture}[
  box/.style={draw=accent, fill=couche, rounded corners=3pt, minimum width=8cm,
              minimum height=1cm, align=center, font=\small},
  side/.style={draw=gray, fill=entite, rounded corners=3pt, minimum width=3.6cm,
               minimum height=1cm, align=center, font=\small},
  arr/.style={-{Stealth[length=2.5mm]}, thick, accent},
  node distance=0.6cm
]
  \node[box] (nav) {Navigateur\\\footnotesize formulaires HTML, liens, barre de recherche};
  \node[box, below=of nav] (vue) {Vues Thymeleaf\\\footnotesize \code{templates/<entite>/\{list,form,show\}.html}, \code{layout/fragments.html}};
  \node[box, below=of vue] (ctl) {Contrôleurs Spring MVC\\\footnotesize 12 contrôleurs CRUD + \code{SearchController} + \code{HomeController}};
  \node[box, below=of ctl] (repo) {Repositories Spring Data\\\footnotesize \code{JpaRepository} + \code{JpaSpecificationExecutor}};
  \node[box, below=of repo] (jpa) {Entités JPA / Hibernate\\\footnotesize 12 entités, \code{ddl-auto: update}};
  \node[box, below=of jpa] (db) {MariaDB\\\footnotesize base \code{hr}};

  \node[side, right=0.7cm of repo] (search) {\code{GenericSearch}\\\footnotesize \code{LIKE} par réflexion};
  \node[side, right=0.7cm of db] (seed) {\code{DataSeeder}\\\footnotesize pays + grades};

  \draw[arr] (nav) -- node[right,font=\footnotesize]{HTTP GET / POST} (vue);
  \draw[arr] (vue) -- (ctl);
  \draw[arr] (ctl) -- (repo);
  \draw[arr] (repo) -- (jpa);
  \draw[arr] (jpa) -- node[right,font=\footnotesize]{JDBC} (db);
  \draw[arr, dashed] (search) -- (repo);
  \draw[arr, dashed] (seed) -- (jpa);
\end{tikzpicture}
\end{center}

\noindent
La flèche entre le navigateur et les vues résume l'aller-retour HTTP : le
contrôleur reçoit la requête, prépare le \code{Model} et renvoie le nom d'un
template que Thymeleaf rend en HTML. \code{open-in-view} est activé, ce qui
permet aux templates de parcourir les collections paresseuses (pages de détail
listant les entités liées) pendant le rendu.

%======================================================================
\section{Organisation du code}

\begin{lstlisting}
src/main/java/fr/lab/hr/
  HrApplication.java   point d'entrée Spring Boot
  model/               12 entités JPA + enum Genre + interface Identifiable
  repository/          12 interfaces Spring Data JPA
  controller/          12 contrôleurs CRUD + SearchController + HomeController
  search/              GenericSearch, SearchHit, SearchResultGroup
  config/              DataSeeder (tables de référence)
src/main/resources/
  application.yml      connexion MariaDB, options JPA et Thymeleaf
  templates/           list/form/show par entité, layout/, search/
\end{lstlisting}

Toutes les entités implémentent \code{Identifiable} (une seule méthode
\code{getId()}), ce qui permet à la recherche de construire des liens vers
n'importe quel résultat sans connaître son type.

%======================================================================
\section{Modèle de données}

\subsection{Entités}

\begin{center}
\begin{tabularx}{\linewidth}{@{}lX@{}}
\toprule
\textbf{Entité} & \textbf{Contenu} \\
\midrule
\code{Personne} & Identité (nom, prénom, prénom d'usage, nom de jeune fille,
  genre), badge, naissance, doctorat, statuts (HDR, associé, collaborateur),
  quotité, ORCID, idHAL, section CNRS, listes de téléphones, e-mails et pages
  web \\
\code{Structure} & Nom, structure parente, sous-structures, bureaux \\
\code{Contrat} & Nom, porteur, sous-porteurs, coordinateurs, partenaires,
  autres personnes \\
\code{Employeur} & Nom, public ou non, grades proposés \\
\code{Grade} & Code et libellé (IE, IR, AI, DR, CR, MCF, PR pré-remplis) \\
\code{Poste} & Personne, employeur, grade, dates de début et de fin \\
\code{Bureau} & Nom, capacité, bâtiment \\
\code{Batiment} & Nom, campus \\
\code{Campus} & Nom \\
\code{Adresse} & Rue, ville, code postal, latitude/longitude, pays \\
\code{Pays} & Code ISO et nom (249 pays pré-remplis) \\
\code{Publication} & Titre, année, identifiant HAL \\
\bottomrule
\end{tabularx}
\end{center}

\subsection{Relations}

La figure~\ref{fig:modele} montre les relations. Chaque flèche part du
\emph{côté propriétaire} (celui qui porte \code{@JoinColumn} ou
\code{@JoinTable}, et donc dont le formulaire édite la relation) et pointe
vers le côté inverse (\code{mappedBy}, lecture seule).

\begin{figure}[htbp]
\centering
\begin{tikzpicture}[
  ent/.style={draw=accent, fill=couche, rounded corners=2pt, minimum width=2.3cm,
              minimum height=0.8cm, font=\small\ttfamily},
  rel/.style={-{Stealth[length=2mm]}, accent},
  lab/.style={font=\scriptsize, fill=white, inner sep=1pt},
  x=4cm, y=2.6cm
]
  \node[ent] (campus)    at (0,2) {Campus};
  \node[ent] (batiment)  at (1,2) {Batiment};
  \node[ent] (bureau)    at (2,2) {Bureau};
  \node[ent] (structure) at (3,2) {Structure};
  \node[ent] (pays)      at (0,1) {Pays};
  \node[ent] (adresse)   at (1,1) {Adresse};
  \node[ent, very thick] (personne) at (2,1) {Personne};
  \node[ent] (publication) at (3,1) {Publication};
  \node[ent] (grade)     at (0,0) {Grade};
  \node[ent] (employeur) at (1,0) {Employeur};
  \node[ent] (poste)     at (2,0) {Poste};
  \node[ent] (contrat)   at (3,0) {Contrat};

  \draw[rel] (batiment) -- node[lab,above]{N:1} (campus);
  \draw[rel] (bureau) -- node[lab,above]{N:1} (batiment);
  \draw[rel] (structure) -- node[lab,above]{N:M} (bureau);
  \draw[rel] (structure.north east) .. controls +(0.6,0.7) and +(-0.6,0.7) .. node[lab,above]{parente N:1} (structure.north west);

  \draw[rel] (adresse) -- node[lab,above]{N:1} (pays);
  \draw[rel] (personne) -- node[lab,above]{N:1} (adresse);
  \draw[rel] (personne) -- node[lab,above]{N:M} (publication);
  \draw[rel] ([xshift=-4mm]personne.north) -- node[lab,left]{principal N:1} ([xshift=-4mm]bureau.south);
  \draw[rel] ([xshift=4mm]personne.north) -- node[lab,right,pos=0.3]{secondaires N:M} ([xshift=4mm]bureau.south);
  \draw[rel] (personne.north east) -- node[lab,sloped,above,pos=0.65]{N:1 et N:M} (structure.south west);

  \draw[rel] (poste) -- node[lab,left]{N:1} (personne);
  \draw[rel] (poste) -- node[lab,above]{N:1} (employeur);
  \draw[rel] (employeur) -- node[lab,above]{N:M} (grade);
  \draw[rel] (poste.south) to[bend left=25] node[lab,below]{N:1} (grade.south);
  \draw[rel] (contrat.north west) -- node[lab,sloped,below,pos=0.45]{porteur N:1 + 4 rôles N:M} (personne.south east);
\end{tikzpicture}
\caption{Modèle de données. Les flèches partent du côté propriétaire de chaque relation.}
\label{fig:modele}
\end{figure}

Toutes les relations sont mappées dans les deux sens : chaque page de détail
affiche et lie les entités associées, quelle que soit la direction.

\subsection{Tables de jointure}

\begin{center}\footnotesize
\begin{tabular}{@{}lll@{}}
\toprule
\textbf{Table} & \textbf{Propriétaire} & \textbf{Vue inverse} \\
\midrule
\code{personne\_bureau\_secondaire} & \code{Personne.bureauxSecondaires} & \code{Bureau.personnesSecondaire} \\
\code{personne\_structure\_impliquee} & \code{Personne.structuresImpliquees} & \code{Structure.personnesImpliquees} \\
\code{personne\_publication} & \code{Personne.publications} & \code{Publication.auteurs} \\
\code{structure\_bureau} & \code{Structure.bureaux} & \code{Bureau.structures} \\
\code{employeur\_grade} & \code{Employeur.grades} & \code{Grade.employeurs} \\
\code{contrat\_sous\_porteur} & \code{Contrat.sousPorteurs} & \code{Personne.contratsSousPortes} \\
\code{contrat\_coordinateur} & \code{Contrat.coordinateurs} & \code{Personne.contratsCoordines} \\
\code{contrat\_partenaire} & \code{Contrat.partenaires} & \code{Personne.contratsPartenaire} \\
\code{contrat\_autre\_personne} & \code{Contrat.autresPersonnes} & \code{Personne.contratsAutre} \\
\bottomrule
\end{tabular}
\end{center}

Les listes de chaînes de \code{Personne} sont des \code{@ElementCollection}
stockées dans \code{personne\_telephone}, \code{personne\_email} et
\code{personne\_page\_web}.

\subsection{Décisions de modélisation}

\begin{description}
  \item[Bureaux d'une personne.] Un bureau principal (\code{ManyToOne}) et une
    liste de bureaux secondaires (\code{ManyToMany}), le cahier des charges
    mentionnant les deux sans les distinguer.
  \item[Identité.] Ajout de \code{nom}/\code{prenom} en plus de
    \code{nomJeuneFille}/\code{prenomUsage}. L'affichage utilise le prénom
    d'usage s'il est renseigné.
  \item[Employeurs et grades.] Plutôt que de sous-classer EPST et Université en
    dur, chaque \code{Employeur} choisit ses \code{Grade} dans un catalogue
    commun. La règle « les employeurs publics définissent des statuts » reste
    ainsi modifiable depuis l'interface.
  \item[Contrats.] Cinq relations distinctes vers \code{Personne}, une par rôle,
    toutes possédées par \code{Contrat}.
  \item[Publications.] Entité minimale (titre, année, idHAL). Les auteurs se
    gèrent depuis la fiche de la personne, côté propriétaire.
\end{description}

%======================================================================
\section{Couche web}

\subsection{Pattern CRUD commun}

Les 12 contrôleurs suivent exactement le même schéma, illustré ici pour
\code{/personnes} :

\begin{center}
\begin{tabular}{@{}lll@{}}
\toprule
\textbf{Méthode} & \textbf{Route} & \textbf{Action} \\
\midrule
GET  & \code{/personnes} & liste (\code{personne/list}) \\
GET  & \code{/personnes/new} & formulaire de création (\code{personne/form}) \\
GET  & \code{/personnes/\{id\}/edit} & formulaire de modification (\code{personne/form}) \\
GET  & \code{/personnes/\{id\}} & page de détail (\code{personne/show}) \\
POST & \code{/personnes} & création ou modification selon le champ caché \code{id} \\
POST & \code{/personnes/\{id\}/delete} & suppression \\
\bottomrule
\end{tabular}
\end{center}

Après un \code{save} ou un \code{delete}, le contrôleur redirige vers la liste
(pattern Post/Redirect/Get). En cas d'erreur de validation, le formulaire est
réaffiché avec les listes de référence nécessaires aux \code{<select>}.

\subsection{Résolution explicite des relations}

Il n'y a pas de \code{Converter<String, Entity>} global. Chaque formulaire
poste les identifiants des entités liées dans des paramètres séparés, que le
contrôleur résout lui-même :

\begin{lstlisting}[language=Java]
@PostMapping
public String save(@Valid @ModelAttribute("personne") Personne personne,
        BindingResult result,
        @RequestParam(required = false) Long bureauPrincipalId,
        @RequestParam(required = false) List<Long> bureauxSecondairesIds,
        ...) {
    personne.setBureauPrincipal(bureauPrincipalId != null
            ? bureauRepository.findById(bureauPrincipalId).orElse(null) : null);
    personne.setBureauxSecondaires(bureauxSecondairesIds != null
            ? bureauRepository.findAllById(bureauxSecondairesIds)
            : Collections.emptyList());
    ...
    repository.save(personne);
    return "redirect:/personnes";
}
\end{lstlisting}

Ce choix rend le comportement de chaque formulaire lisible sans connaître la
configuration globale de Spring, au prix d'un peu de répétition.

\subsection{Listes de chaînes}

Les listes simples (téléphones, e-mails, pages web) sont éditées dans un
\code{<textarea>}, une valeur par ligne, et découpées côté serveur par
\code{linesToList}. Aucun JavaScript n'est nécessaire.

\subsection{Vues}

Chaque entité possède trois templates : \code{list.html} (tableau, liens
d'ajout et de suppression), \code{form.html} (création et modification,
\code{<select>} simples ou multiples pour les relations) et \code{show.html}
(détail, avec les entités liées cliquables dans les deux sens). La barre de
navigation et la barre de recherche sont factorisées dans
\code{layout/fragments.html}.

%======================================================================
\section{Recherche globale}

Aucune requête de recherche n'est écrite à la main. \code{GenericSearch}
construit par réflexion une \code{Specification} JPA qui applique un
\code{LIKE} insensible à la casse sur tous les champs \code{String} déclarés de
l'entité :

\begin{lstlisting}[language=Java]
public static <T> Specification<T> containing(Class<T> entityClass, String term) {
    String pattern = "%" + term.toLowerCase() + "%";
    return (root, query, cb) -> {
        List<Predicate> predicates = new ArrayList<>();
        for (Field field : entityClass.getDeclaredFields()) {
            if (field.getType().equals(String.class)) {
                predicates.add(cb.like(cb.lower(root.get(field.getName())), pattern));
            }
        }
        if (predicates.isEmpty()) {
            return cb.disjunction();
        }
        return cb.or(predicates.toArray(new Predicate[0]));
    };
}
\end{lstlisting}

\code{SearchController} (\code{GET /recherche?q=...}) applique cette
spécification aux 12 repositories, regroupe les résultats par type
(\code{SearchResultGroup} de \code{SearchHit}) et masque les groupes vides.
Une nouvelle entité devient cherchable en l'ajoutant simplement à la liste de
\code{SearchController.rechercher()}.

La recherche ne porte que sur les champs texte directs : chercher le nom d'une
structure fera remonter la structure, pas les personnes qui y sont rattachées.

%======================================================================
\section{Persistance et configuration}

\begin{itemize}
  \item \textbf{Schéma.} \code{spring.jpa.hibernate.ddl-auto: update} :
    Hibernate crée et fait évoluer le schéma au démarrage. Il n'y a pas de
    migrations Flyway/Liquibase.
  \item \textbf{Base locale.} MariaDB installée sur le poste (Homebrew),
    port 3306, base \code{hr}, utilisateur \code{hr}/\code{hr}, créés par
    \code{scripts/creer\_base\_dev.sql}.
  \item \textbf{Données de référence.} \code{DataSeeder}
    (\code{ApplicationRunner}) remplit \code{pays} à partir de
    \code{Locale.getISOCountries()} (noms en français) et \code{grade} avec
    IE, IR, AI, DR, CR, MCF et PR. Il ne fait rien si les tables contiennent
    déjà des données.
  \item \textbf{Développement.} Cache Thymeleaf désactivé et
    \code{spring-boot-devtools} pour le rechargement à chaud.
\end{itemize}

\begin{lstlisting}
brew services start mariadb   # MariaDB locale
mvn spring-boot:run      # application sur http://localhost:8080
\end{lstlisting}

%======================================================================
\section{Invariants et pièges connus}

\subsection{Pas de cascade sur les collections non resoumises}

Chaque \code{POST} reconstruit l'entité à partir du formulaire via
\code{@ModelAttribute}. Les collections absentes du formulaire reviennent donc
à leur valeur par défaut, en général une liste vide. Si une telle collection
porte \code{cascade = CascadeType.ALL, orphanRemoval = true}, Hibernate
supprime silencieusement toutes les lignes enfants à la sauvegarde suivante du
parent.

C'est arrivé avec \code{Personne.postes} : chaque enregistrement d'une
personne effaçait ses postes. La correction a été de retirer cascade et
orphanRemoval ; \code{postes} est désormais une vue inverse en lecture seule,
les postes étant gérés par leur propre contrôleur.

\textbf{Règle :} si l'entité enfant a son propre contrôleur CRUD, ne pas
mettre de cascade côté parent.

\subsection{Les côtés inverses ne s'éditent pas}

Les collections \code{mappedBy} (par exemple \code{Bureau.personnesPrincipal}
ou \code{Publication.auteurs}) sont uniquement affichées. Une relation se
modifie toujours depuis le formulaire de l'entité propriétaire.

\subsection{Toute nouvelle relation est bidirectionnelle}

Le cahier des charges exige la navigation dans les deux sens : toute relation
ajoutée doit être mappée des deux côtés, éditée via des paramètres
\code{xxxId}/\code{xxxIds} côté propriétaire, et affichée sur les deux pages de
détail.

%======================================================================
\section{Limites et évolutions possibles}

\begin{itemize}
  \item Pas d'authentification ni de gestion des droits (voulu pour la v1).
  \item \code{ddl-auto: update} ne convient pas à la production : passer à des
    migrations versionnées (Flyway) avant tout déploiement réel.
  \item Le formulaire \code{Poste} propose tous les grades, sans les filtrer
    selon ceux de l'employeur choisi.
  \item La recherche ne traverse pas les relations.
  \item Les tests automatisés se limitent au test de chargement du contexte
    (\code{HrApplicationTests}) ; toute règle métier non triviale ajoutée
    devrait être accompagnée d'un test.
\end{itemize}

\end{document}
